\scanpage{42}
Moyennant $C(Y)$, $C(X)$ équivaut à la condition :

\noindent\begin{tabular}{@{}|p{.92\linewidth}@{}}
$\forall i\leq n-1$, l'hom $\psi_\xi^i\colon H^i(Y)\longrightarrow H^i(X)$, inverse à gauche naturel de $\varphi^i\colon H^i(X)\unsure{\mathord{-}|\longrightarrow}H^i(Y)$, \deleted{\unsure{est}} égal à $\Lambda_X\varphi_{\unsure{i|*}}\mid H^i(Y)$, est induit par une classe algébrique.
\end{tabular}

ou encore, :

\noindent\begin{tabular}{@{}|p{.92\linewidth}@{}}
$\Lambda_X\varphi_*$ est algébrique
\end{tabular}

\deleted{ou encore} N.B. $\Lambda_X\varphi_*$ correspon\deleted{\unsure{d à|d}}d à une classe \inserted{dans $\sum H^{2(n-1)-i}(Y)\otimes H^i(X)=H^{2(n-1)}(X\times Y)=H^{m-1}(X\times Y)$, $m=\dim X\times Y$} de cohomologie sur $X\times Y$, que nous polarisons par $\xi_X\otimes1+1\otimes\xi_Y$. Or \deleted{\unsure{$\xi_X\otimes1+$|$\zeta_X\otimes1+$}} $L_{X\times Y}(\Lambda_X\varphi_*)$ correspond à
\[
\begin{aligned}
\underbrace{\Lambda_X\varphi_*L_Y}_{\substack{\deleted{\unsure{\Lambda_X\varphi_*\varphi^*}}\\ \Lambda_XL_X\varphi_*}}
+\underbrace{L_X\Lambda_X\varphi_*}_{\deleted{\unsure{\varphi_*\Lambda_Y\varphi^*\varphi_*}}}
&=\deleted{\left(2\,\mathrm{id}-\sum_{n\leq i\leq2n}p_i^X\right)\varphi_*}\\
&(\Lambda_X L_X+L_X\Lambda_X)\varphi_*\\
&=(\varphi_*\Lambda_Y\varphi^*+\mathrm{id}_X)\varphi_*
\end{aligned}
\]
(car $p_n^X\varphi_*=0$)

Donc c'est algébrique. Donc on trouve
\[
 C(X)\Longleftrightarrow\bigl[C(Y)\text{ et }A^\circ(X\times Y)\bigr]
\]
puis de proche en proche
\[
 C(X)\Longleftrightarrow\bigl(A^\circ(X\times Y)\text{ et }A^\circ(Y\times Z)\text{ et }\deleted{\unsure{A}}A^\circ(Z\times T)\;{-}{-}\bigr)
\]
\deleted{$X\times Y,\,Y\times\unsure{Z}$} $X\supset Y\supset Z\supset T\,{-}{-}$ est une suite décroissante de sections hyperplanes non singulières.

% stamp: repeated diagonal photographic-paper watermark, apparently Gevaert.
% archive sheet number: 41

